The first module of the course focuses on \emph{defining} a language model mathematically.
To see why such a definition is nuanced, we are going to give an informal definition of a language model and demonstrate two ways in which that definition breaks and fails to meet our desired criteria.
\begin{definition}[Informal]\label{def:language-model-informal}
Given an alphabet\footnote{An alphabet is a \emph{finite}, non-empty set.} $\alphabet$ and a distinguished end-of-sequence symbol $\eos \not\in \alphabet$, a language model is a collection of conditional probability distributions $\pdens(\sym \mid \str)$ for $\sym \in \alphabet \cup \{\eos\}$ and $\str \in \kleene{\alphabet}$, the Kleene closure of $\alphabet$.\footnote{We will define the Kleene closure in more detail later in the course.}
$\pdens(\sym \mid \str)$ therefore represents the probability of $\sym$ being the next token given the history $\str$.
\end{definition}
\Cref{def:language-model-informal} is the definition of a language model that is implicitly assumed in most papers on language modeling.
We say implicitly since most technical papers on language modeling simply write down the following autoregressive factorization
\begin{equation}\label{def:autoregressive-lm}
    \pdens(\str) = \pdens(\sym_1\cdots \sym_\strlen) = \pdens(\eos \mid \str_{\leq \strlen}) \prod_{\tstep=1}^\strlen \pdens(\sym_t \mid \str_{<\tstep})
\end{equation}
as the probability of a string according to the distribution $\pdens$.\footnote{Many (erroneously) authors avoid writing $\eos$ for concision.}
The part that is left implicit in \cref{def:autoregressive-lm} is whether or not $\pdens$ is indeed a probability distribution and, if it is, over what space.
The natural assumption in \cref{def:language-model-informal} is that $\pdens$ is a distribution over $\kleene{\alphabet}$, i.e., the set of all \emph{finite} strings\footnote{Some authors assert that strings are by definition finite.} over an alphabet $\alphabet$.
However, in general, it is not true that all such collections of conditionals will yield a valid probability distribution over $\kleene{\alphabet}$; some may ``leak'' probability mass to infinite sequences.\footnote{However, the converse \emph{is} true: All valid distributions over $\kleene{\alphabet}$ may be factorized as the above.}
More subtly, we additionally have to be very careful when defining \emph{which} events $\pdens$ assigns a probability to lest we run into a classic paradox, which we will demonstrate below.\looseness=-1

We highlight the two issues above with two very simple examples.
The first example is a well-known paradox in probability theory.
\begin{example}[Infinite Coin Toss]\label{ex:inf-coin-toss}
Consider the infinite independent fair coin toss model, where we aim to place a distribution over $\{\texttt{H},\texttt{T}\}^\infty$, the (uncountable) set of infinite sequences of $\{\texttt{H},\texttt{T}\}$ (\texttt{H} represents the event of throwing heads and \texttt{T} the event of throwing tails). 
Intuitively, such a distribution corresponds to an ``language model'' as defined above in which for all $\str_{<\idx}$, $\pdens(\texttt{H} \mid \str_{<\idx}) = \pdens(\texttt{T} \mid \str_{<\idx}) = \frac{1}{2}$ and $\pdens(\eos \mid \str_{<\idx})=0$.
However, each individual infinite sequence over $\{\texttt{H},\texttt{T}\}$ should also be assigned probability $(\frac{1}{2})^\infty = 0$. 
Without a formal foundation, arrives at the following paradox:\looseness=-1
\begin{align*}
    1&=\pdens\left(\{\texttt{H},\texttt{T}\}^\infty\right) 
    =\pdens\left(\bigcup_{\bomega\in\{\texttt{H},\texttt{T}\}^\infty} \{\bomega\}\right) \\
    &=\sum_{\bomega\in\{\texttt{H},\texttt{T}\}^\infty}\pdens(\{\bomega\}) = \sum_{\bomega\in\{\texttt{H},\texttt{T}\}^\infty} 0 \stackrel{?}{=} 0. \nonumber
\end{align*}
\end{example}
The second example is more specific to language modeling.
As we stated above, an implicit assumption made by most language modeling papers is that a language model constitutes a distribution over $\kleene{\alphabet}$. 
However, in our next example, we show that a collection of conditions that satisfy \cref{def:language-model-informal} may not sum to 1 if the sum is restricted to elements of $\kleene{\alphabet}$. 
This means that it is not a-priori clear what space our probability distribution is defined over.\footnote{This also holds true for the first example.}
\begin{example}[Possibly Finite Coin Toss]\label{ex:possibly-inf-coin-toss}
Consider now the possibly finite ``coin toss'' model with a rather peculiar coin: when tossing the coin for the first time, both \texttt{H} and \texttt{T} are equally likely.
After the first toss, however, the coin gets stuck: If $\sym_1 = \texttt{H}$, we can only ever toss another \texttt{H} again, whereas if $\sym_1 = \texttt{T}$, the next toss can result in another \texttt{T} or ``end'' the sequence of throws ($\eos$) with equal probability. 
Formally, we model a probability distribution over $\kleene{\left\{\texttt{H}, \texttt{T}\right\}} \cup \left\{\texttt{H}, \texttt{T}\right\}^\infty$, the set of finite and infinite sequences of tosses.\ryan{There surely is a better way to format this.} 
Formally, 
$\pdens(\texttt{H} \mid \varepsilon) = \pdens(\texttt{H} \mid \str_{<1}) = \pdens(\texttt{T} \mid \varepsilon) = \pdens(\texttt{T} \mid \str_{<1}) = \frac{1}{2}$, 
$\pdens(\texttt{H} \mid \str_{<\idx}) = 
\begin{cases}
    1 & \text{ if } \sym_{\idx - 1} = \texttt{H} \\ 
    0 & \text{ otherwise}
\end{cases}$, 
$\pdens(\texttt{T} \mid \str_{<\idx}) = 
\begin{cases}
    \frac{1}{2} & \text{ if } \sym_{\idx - 1} = \texttt{T} \\ 
    0 & \text{ otherwise}
\end{cases}$ 
for $\idx > 1$, and $\pdens(\eos \mid \str_{<\idx})= \begin{cases}
    \frac{1}{2} & \text{ if } \idx > 0 \text{ and } \sym_{\idx - 1} = \texttt{T} \\
    0 & \text{ otherwise.}
\end{cases}$.

If you are familiar with (weighted) finite-state automata,\footnote{They will be formally introduced in \cref{sec:n-gram-models}} you can imagine the model as depicted in \cref{fig:possibly-finite-coin-toss}.
\begin{figure}[ht]
    \centering
    \footnotesize
    \begin{tikzpicture}[minimum size=7mm]
        \node[state, initial] (q0) {$0/1$};
        \node[state,above right=0.5cm and 0.8cm of q0] (q1) {$1$};
        \node[state,below right=0.5cm and 0.8cm of q0,accepting] (q2) {$2/\frac{1}{2}$};
        
        \path (q0) edge [bend left] node [above, sloped] {$\texttt{H}/\frac{1}{2}$} (q1);
        \path (q0) edge [bend right] node [below, sloped] {$\texttt{T}/\frac{1}{2}$} (q2);
        \path (q1) edge [loop right] node [right] {$\texttt{H}/1$} (q1);
        \path (q2) edge [loop right] node [right] {$\texttt{T}/\frac{1}{2}$} (q2);
    \end{tikzpicture}
    \caption{Graphical depiction of the possibly finite coin toss model.}
    \label{fig:possibly-finite-coin-toss}
\end{figure}

It is easy to see that this model only places the probability of $\frac{1}{2}$ on \emph{finite} sequences of tosses.
If we were only interested in those (analogously to how we are only interested in finite strings when modeling language), yet still allowed the model to specify the probabilities as in this example, the resulting probability distribution would not model what we require.
\end{example}

It takes some mathematical heft to define a language model in a manner that avoids such paradoxes. 
The tool of choice for mathematicians is measure theory, as it allows us to define
probability over uncountable sets\footnote{As stated earlier, $\{\texttt{H},\texttt{T}\}^\infty$ is uncountable. It's easy to see there exists a bijection between $\{\texttt{H},\texttt{T}\}^\infty$ and the binary expansion of the real interval $[0, 1]$.} in a principled way.
Thus, we begin our formal treatment of language modeling with a primer of measure theory in \cref{sec:measure-theory}.
Then, we will use concepts discussed in the primer to work up to a formal definition of a language model.

% \ryan{Move up the paradox example to here.}
% \anej{Show that you can get the probability of $0.5$ on finite strings and $0.5$ on infinite ones.}